% In-paper appendix (MVA, single PDF): records and computations the referee needs to see. The rest: Online Resource 1 (esm.tex).
\section{Experiment log summary and computation notes}\label{sec:deney-kaydi}

This appendix gives the records and computations that a referee needs to check. It covers a summary of the run log, the application of the stopping rule, post-unlock work and the GPU budget. It also covers the computation of the screening probabilities, the gradient-norm ratio and the sources of the plausibility anchor. The run-log table, additional results and method details are in Online Resource 1 (also in the GitHub repository).

Table~\ref{tab:deney-kaydi} in Online Resource 1 breaks down the runs per card. It covers the baseline plausibility check, the three rounds, the locked-test stage and the post-test additional-seed runs. Of the 60 run stamps in total, five did not end with exit code 0. Three are calibration runs (rc=98), one is an orphan process (rc=96) and one is an H05 failure (rc=95). There are 35 runs of evidential value. The first 22 are 12 stage-3 trainings, 1 H07 evaluation and 9 locked-test evaluations. The other 13 are 12 trainings for additional seeds and 1 WBF diagnostic run. The remaining 25 runs are gate, control, calibration and failure records and were not reported as results. A total of 75.45 GPU hours was spent (target 75, ceiling 90). The full command and rule parameters of each run are stored in the \texttt{config.json} file in the run directory. The queues are stored under \texttt{kuyruk/}.

\subsection{Application of the stopping rule}\label{sec:ek-durma}
Fig.~\ref{fig:f4-butce}a in Online Resource 1 shows the GPU time per round and the cumulative total. Panel (b) shows the $\Delta$ of the top training-time card in each round. Filled circles are three-seed means, and open circles are seed-0 screening values. The stopping rule says: move to the locked test if the top training card has $\Delta<+0.30$ in two consecutive rounds. It was triggered at the end of the second round, based on the values available at that date. The top training card of the first round was H03, and its three-seed mean (+0.25) was already known. The second round had only seed-0 screening runs, and its top card was H04 ($-$0.30). The cumulative time at that point was 35.30 hours. The second- and third-round means in the figure were computed with post-unlock additional seeds. They did not enter the stopping decision.

The rule text does not specify whether $\Delta$ is the seed-0 value or the three-seed mean. For each round, we used the best estimate available at the time. This was the three-seed mean in the first round (+0.25) and the seed-0 value in the second ($-$0.30). Had both rounds been read with the seed-0 value, the first-round value would have been the seed-0 difference of H03. That value is +0.37, i.e., $\ge+0.30$, so the rule would not have been triggered at this point. The rule text also does not distinguish training-time cards from training-free ones. In practice, the training-free TTA row of the second round was kept outside the rule. This row was on validation with seed 0, gave +0.80, and fell outside the $\times$1.10 relative latency constraint. Had TTA also counted as a round-2 card, the rule would not have been triggered (+0.80 $\ge$ +0.30). The stopping decision therefore rests on the training-time interpretation.

H08 was the only card planned for the dose question. It was run in the third round, after the rule had been triggered. Its screening value was +0.22 (3.16 hours), and its three-seed mean was $-$0.08. The study then moved to the locked test (0.21 hours). The design, code and run expectation of H08 had been written about 9.6 hours before the trigger, before any second-round result had arrived. The execution sequence was established from the queue and code file logs. The rule document set a limit of five rounds and contained a candidate pool for the third round. However, it had no pre-specified exception that allowed another round after the rule had been triggered. Running H08 is therefore a deviation from the pre-specified rules (Section~\ref{sec:sinirliliklar}). The card was not accepted, and the deviation did not turn any decision positive.

\subsection{Post-unlock work}\label{sec:ek-kilit}
The following is the complete list of work done on VOC after unlocking.
(i) Seed 1 and seed 2 validation runs of the six cards eliminated at screening (H01, H02, H04, H05, H06, H08), i.e., 12 training runs.
(ii) An exploratory diagnostic grid for the WBF branch of H07 on the validation set. It varied the WBF threshold, the skip threshold and the fusion type. None of its settings was applied to the test, and the best point of the grid does not count as evidence.
(iii) Analyses of n=3, power and chance. These used existing outputs and required no new runs.
(iv) The two component cards of H05 (H05-pos, H05-neg), run under a second rule document (Section~\ref{sec:h05-mekanizma}). They were run with three seeds on validation only (6 training runs), and the VOC test was not opened.
(v) H05c, the $\lambda_{\mathrm{cls}}$ gain compensation card for H05, run under a third rule document (2026-10-22). It was run with three seeds on validation only (3 training runs), and the VOC test was not opened.

This list covers work on VOC data. The KITTI experiment was carried out on a separate dataset under a separate rule document (Section~\ref{sec:genelleme}). After the KITTI test was opened, no new training, test inference or hyperparameter search was done. The pre-specified parts were the test evaluation of the 48 runs, the H07 test rows and the primary analysis. Everything else is exploratory or secondary. This includes the pre-specified secondary analysis on validation and the re-evaluation of \texttt{best.pt} and \texttt{last.pt} (96 evaluations). It also includes several analyses of the existing test predictions: clustered bootstrap $p$, Holm-10, Holm-14 and sign-flip sensitivity analyses, $\Delta$ without \texttt{DontCare}, class subsets and secondary tables. The remaining items are the hierarchical bootstrap, $\Delta$ per scene group, the tram sensitivity analysis, the external reference value and a diagnosis of the baseline on 1000 training images (Online Resource 1, Section~\ref{sec:ek-ikincil}). Apart from these, no other runs, hypotheses or hyperparameter searches were carried out on VOC, and the test was not touched. The additional seeds were added post-unlock, following the principle of not reporting numbers with n<3. They were given only to the six cards eliminated by the pre-unlock screening decisions.

\subsection{GPU budget}\label{sec:ek-butce}
After unlocking, the additional-seed runs of the six cards and the WBF diagnostic grid took 36.79 hours. They raised the total from 38.67 to 75.45 hours; totals were computed from unrounded values. This is slightly above the target (75) and below the ceiling (90).

The generalization experiment (Section~\ref{sec:genelleme}) is not included in Table~\ref{tab:deney-kaydi} (Online Resource 1) and was logged separately. It comprises 54 full runs, all with 100 epochs. On KITTI, two models $\times$ eight arms $\times$ three seeds give 48 runs. On VOC, H05-pos and H05-neg $\times$ three seeds give 6 runs. The total training time is 89.02 hours: YOLO11s 41.15, YOLO26n 28.99 and the VOC mechanism cards 18.88. The three runs of H05c (VOC, validation; 10.25 GPU hours) were logged separately and are not included in this total. Smoke tests, short pilots, H07 and test evaluation are also excluded. This total is therefore defined differently from the 75.45-hour total of the first experiment, which counts all run stamps. The KITTI test split was opened on 2026-10-22 at 12:02:15. It was opened in a single opening with the \texttt{best.pt} and \texttt{last.pt} weights of the 48 runs. The VOC test was not opened in this experiment.

\subsection{Computation of the screening probabilities}\label{sec:ek-onele}
The two probabilities in Section~\ref{sec:bulgu-ablasyon} answer different questions and are descriptive only.

(A) The first is the probability that an uninformative seed-0 gate ($\Delta_{k,0}\ge+0.30$) eliminates all six cards when their true effect is zero. It was computed as $(1-p)^6$, where $p=1-\Phi(0.30/\sigma)$ is the probability that a single card passes the gate. A single common $\sigma$ is used, with two candidate values. The first is an image-sampling estimate, obtained by scaling the test bootstrap standard error of H03 to the validation size (0.502; $p=0.275$). The second is the s.d. of the validation seed differences of H03 (0.240; $p=0.105$). These give 0.145 and 0.512, respectively.

(B) The second probability uses the definition ``an eliminated card is not accepted with three seeds either''. Its value, 0.774, was obtained by Monte Carlo with 400\,000 draws per card. Each draw has three seeds with $\Delta_s\sim N(0,s_{\Delta,k}^2)$, where $s_{\Delta,k}$ is the observed s.d. of the card's seed differences. The true effect is exactly zero. A card is accepted if its mean $\Delta\ge+0.40$ and the one-sided $t>1.886$. Cards were drawn independently. Two dependencies were not modeled: the correlation arising from the common baseline, and the inclusion of the screening seed in the $n=3$ mean.

Neither computation proves that the gate is or is not informative. The sensitivity of the gate cannot be measured, because no true positive was observed.

\subsection{Gradient-norm ratio}\label{sec:ek-gradyan}
The VFL/BCE ratio reported for H05 is computed from the L2 norm of the loss gradient with respect to the classification logits. For the same logits and targets, the VFL and BCE gradient norms are measured every 200 training steps, giving 486 measurements per run. The ratio of the norm sums is then taken. The effective gradient ratio in the manipulation check of H05c is this ratio multiplied by $1.47/0.5$ (Section~\ref{sec:h05}).

\subsection{Sources of the plausibility anchor}\label{sec:ek-cipa}
Two YOLO11n sources give the lower end of the anchor in Section~\ref{sec:replikasyon}. ADL-YOLO reports training from scratch, whereas YOLO-Master does not specify the initialization \cite{zhang_2025_adlyolo,lin_2025_yolomaster}. The YOLOv8n value at the upper end (63.5) was verified in the text of Section 4.4 of the source. The corresponding table cell could not be read. The source uses the VOC07 test set as its validation or test set, so its checkpoint was probably selected on the test set \cite{elhamdouni_2025_yolov8nvoc}. The value of 60.2 for YOLOv8n trained from scratch lies outside the anchor and is listed for context only (Table~\ref{tab:replikasyon}).
